\documentclass[main.tex]{subfiles}
\begin{document}

\section{Introduction*}
\label{sec:intronew}

Modern computing is fundamentally an interactive experience.
And for most users, a screen is the primary interface to a computer.
It is through screen through which people browse webpages, interact with the
desktops or play video games.
Although these applications describe visual information in different ways they
must be displayed as pixels on a screen.
The process of converting this high-level visual information into a collection
of pixels is called rasterization.
While some applications, like video games, have to raster three-dimensional
scenes, others, like browsers, need to render only flat shapes, images and text.
Such applications are powered by 2D rasterization libraries like Skia,
CoreGraphics, and Direct2D.
While these applications lack the complexity introduced by the third dimension,
  they still need to render high level descriptions of many shapes with
  complex effects, gradients, and sometimes even
  animations.
Moreover,
  they must do so at interactive rates (like 60 to 120 frames per second) to
  ensure a smooth experience for end users.
Hence, modern 2D rasterization libraries make the best use of modern hardware
  as possible.
For example, Skia, the 2D rasterization library of choice of most browser
  engines, supports multiple backends, including GPU accelerated ones.
Its newest GPU backend, Graphite, is designed to take better advantage of modern
  Graphics APIs like Metal and Vulkan.
However, a fast rasterizeration library is not enough to ensure smooth
experience.
\fillin{i want to say rendering or viewing as an adjective dunno tho}.
Clients of these libraries still need to make effective use of them.
But this is difficult for even sophisticated and well-engineered clients.
Google Chrome is developed alongside Skia at Google, yet it produces
\fillin{highly} suboptimal Skia \fillin{programs. We never introduce the notion
  that these clients are taking programs as inputs.}
The main reason for this suboptimiality is the dsiconnect between the data
models of the client and the library.
Chrome describes visual information as HTML and CSS documents,
  while Skia needs
  a sequence of drawing commands describing what shapes to draw, how to draw
  them, and how to blend them together.
Moreover, not all skia commands have the same performance characteristics.
Some may allocate extra pixel buffers or perform additional blends.
Different sequences of Skia commands may draw the same image, but may take
drastically different amounts of time to draw.
It is upto the client to generate the \emph{right} sequence of commands.
But how is a client supposed to do this?

The relationship between the client and the library is quite
  unlike that of a compiler of a high-level-language and assembly.
The compiler needs to reason about assembly instructions in order
to efficiently compile programs written in the high-level language.
The msot straightforward way to do this, would be to write doen the semantics of
the assembly instructions we are targetting.

In this paper we aim to do something similiar for rasterization libraries.
Specifically, we introduce \muskia, a formalization of the command language of
the Skia library.
 it captures key features
  like canvas state, the layer stack, clipping, and blending.
To facilitate automated reasoning and verification,
  \muskia has a layered structure.
The top layer is a command language with an operational semantics.
This operational semantics is phrased in terms of a functional core,
  which forms the middle layer.
This functional core then denotes
  to an abstract model of 2D graphics, which is the bottom layer.
The abstract model permits multiple instantiations
  to allow reasoning modulo concerns like anti-aliasing and discretization error.
The semantics including various extensions, is mechanized in Lean.

We then use the \muskia semantics to develop
  an effective and efficient optimizer for Skia.
We identify 4 rewrite rules
  that optimize common Skia command patterns
  generated by Google Chrome on real-world websites.
We verify each rewrite rule using the \muskia semantics,
  and then develop an optimizer that applies these rewrite rules.
The optimizer is highly efficient,
  optimizing even large Skia programs in microseconds,
  and significantly reduces rasterization time.
On a collection of \grmtlnumbench Chrome-generated Skia programs
  from the top 100 websites,
  the optimizer is able to improve rasterization time on \grmtlnumspmatch,
  with speedups ranging from \grmtlspmatchmin to \grmtlspmatchmax.
Moreover,
  the optimizer produces optimization traces
  that can be loaded into Lean
  for translation validation against our semantics.
The optimizer is effective across Skia backends and GPUs,
  and is also effective on Skia programs not seen by the authors.

\medskip
\noindent
The contributions to this paper are:
\begin{enumerate}
\item A new formalization of the Skia API (\Cref{sec:semantics})
\item A collection of rewrites proven valid using this formalization (\Cref{sec:rewrites})
\item A optimizer that speeds up real world Skia programs using these rewrites (\Cref{sec:compiler}).
\end{enumerate}






\section{Introduction}
\label{sec:intro}

Browser engines are an increasingly prominent way
  that users interact with applications---%
  not only web applications like Office~365 or GMail
  but also desktop applications like Slack, Discord, and VS~Code.
Browser applications are graphical,
  and increasingly support complex graphical effects
  such as shadows, gradients, transparency, and animations.
Yet drawing these visual effects 60 or even 120 times per second,
  to keep up with smooth animations or rapid user interaction,
  is challenging.
Fast rasterization,
  the task of actually determining the color of each pixel
  drawn by the application, is critically.
Typically, a rasterization library
  takes as input a high-level description
  of shapes to be drawn and effects to be applied,
  and must then color in millions of pixels
  in the span of hundreds of microseconds.

This requires highly-optimized rasterization libraries.
The most prominent of these libraries is Skia.
Skia is used by nearly all browser engines---%
  it is the only rasterization library
  for the Google Chrome, Mozilla Firefox, and Ladybird browsers,
  and the Linux rasterization library
  for the Webkit engine behind Apple Safari.
Skia is additionally used in the Android operating system,
  and is one of two rasterization options in Flutter,
  a popular mobile application framework.
Skia is also a popular backend for various libraries
  for rastering PDFs, images, and icons,
  and it is also used by a variety of other native applications,
  like for example Sublime Text.
Skia clients describe shapes to be drawn,
  define how different shapes are blended together,
  and apply effects to the results.
Skia then converts this description
  into an efficient sequence of GPU shader calls
  using GPU backends tuned for various native APIs
  like Metal, Vulkan, and OpenGL.

To achieve max performance,
  clients must use Skia the right way.
Certain Skia commands, for example,
  allocate additional pixel buffers,
  perform additional blend operations,
  or switch between special-case and general-purpose masking routines.
Two different sequences of commands might draw the same image
  but with radically different rasterization time.
Even sophisticated clients like Google Chrome
  generate subpar command sequences
  due to the complexity of the Skia command language
  and the clients' own internal complexity.
Ideally, these concerns would be handled internally,
  by an optimization pass within the Skia library,
  but to date no such optimization pass exists.

We address this with
  a new formal semantics of Skia, called \muskia,
  and leverage it to develop such an optimization pass.
\muskia is a formalization of the command language
  implicit in the Skia API; it captures key features
  like canvas state, the layer stack, clipping, and blending.
To facilitate automated reasoning and verification,
  \muskia has a layered structure.
The top layer is a command language with an operational semantics.
This operational semantics is phrased in terms of a functional core,
  which forms the middle layer.
This functional core then denotes
  to an abstract model of 2D graphics, which is the bottom layer.
The abstract model permits multiple instantiations
  to allow reasoning modulo concerns like anti-aliasing and discretization error.
The semantics including various extensions, is mechanized in Lean.

We then use the \muskia semantics to develop
  an effective and efficient optimizer for Skia.
We identify 4 rewrite rules
  that optimize common Skia command patterns
  generated by Google Chrome on real-world websites.
We verify each rewrite rule using the \muskia semantics,
  and then develop an optimizer that applies these rewrite rules.
The optimizer is highly efficient,
  optimizing even large Skia programs in microseconds,
  and significantly reduces rasterization time.
On a collection of \grmtlnumbench Chrome-generated Skia programs
  from the top 100 websites,
  the optimizer is able to improve rasterization time on \grmtlnumspmatch,
  with speedups ranging from \grmtlspmatchmin to \grmtlspmatchmax.
Moreover,
  the optimizer produces optimization traces
  that can be loaded into Lean
  for translation validation against our semantics.
The optimizer is effective across Skia backends and GPUs,
  and is also effective on Skia programs not seen by the authors.

\medskip
\noindent
The contributions to this paper are:
\begin{enumerate}
\item A new formalization of the Skia API (\Cref{sec:semantics})
\item A collection of rewrites proven valid using this formalization (\Cref{sec:rewrites})
\item A optimizer that speeds up real world Skia programs using these rewrites (\Cref{sec:compiler}).
\end{enumerate}
\end{document}
